\section{Introduction}

Introduce BDI agents as a framework in which plan libraries provide the reusable procedural knowledge for achieving goals under different belief contexts.

Explain that BDI plan libraries are traditionally handcrafted by developers who must encode domain behaviour and anticipate applicable contexts, creating a substantial knowledge-engineering burden when adapting agents to new domains.

Relate this engineering burden to the practical difficulty of deploying BDI agents when a sufficiently complete and reliable plan library is not already available.

Distinguish a grounded plan for one planning instance from a reusable BDI plan template that must apply across different objects, states, and future goal invocations.

Explain that generating reusable plan templates is harder than solving one planning instance because applicability conditions, parameters, subsidiary-goal structure, and correctness must remain valid beyond the state in which the plan was generated.

Review existing BDI-planning approaches that invoke planners at runtime or use planning during means--ends reasoning to obtain behaviour for the current goal.

Identify reusable plan-library generation as a different problem because the generated behaviour must remain available inside the BDI agent and support future goal invocations.

Motivate generalized planning as a promising source of reusable behaviour because it generalizes across families of related planning instances rather than solving only one grounded problem.

Identify the remaining gap between generalized-planning behaviour and an executable BDI plan library because the available behaviour may not provide all reusable subsidiary plans required by the agent.

Introduce achievement goals and temporally extended goals as the two classes of user queries supported by the generated plan library.

Introduce controlled natural-language queries as the user-facing way to invoke the generated library for both achievement and temporally extended goals.

Introduce GP2PL as a framework that first transforms generalized-planning behaviour and domain knowledge into a reusable BDI plan library and then reuses that library to handle subsequent user queries.

Present the generation of a reusable and validated BDI plan library from generalized-planning behaviour as the first contribution.

Present a natural-language query interface that reuses the same generated library for both achievement and temporally extended goals as the second contribution.
